% !TeX program = lualatex
% !TeX encoding = utf8
% !TeX spellcheck = uk_UA
% !TeX root = ../main.tex

%% --------------------------------------------------------
%% Chapter file: Confinement.tex
%% --------------------------------------------------------

\graphicspath{{\currfilebase/Pictures}}

\Opensolutionfile{answer}[\currfilebase/\currfilebase-Answers]
\Writetofile{answer}{\protect\section*{Розділ~\ref*{sec:magnetic-confinement}. \nameref*{sec:magnetic-confinement}}}
\chapter{Магнітне утримання плазми}\label{sec:magnetic-confinement}\hypertarget{magnetic-confinement}{}
%% --------------------------------------------------------

Щоб отримати термоядерні реакції, плазму потрібно нагріти до високої температури
й~водночас не дати їй торкнутися стінок. Магнітне поле розв'язує цю задачу для заряджених
частинок лише частково: воно сильно обмежує рух поперек поля, але саме по собі не
перешкоджає руху вздовж силової лінії. Тому геометрія магнітного поля принципово важлива.

\section{Проблема утримання}
%% --------------------------------------------------------

Для реакції синтезу $\mathrm D+\mathrm T\to{}^4\mathrm{He}+n$ дейтерієво-тритієву плазму
треба нагріти до температури $T\sim10$--$20$~кеВ ($\sim10^8$~К) і~утримувати
достатньо довго. Критерій Лоусона (розд.~\ref{sec:fusion}) вимагає
$n\tau_E\gtrsim10^{14}$--$10^{15}~\mathrm{см}^{-3}\,\mathrm{с}$ (оцінка за порядком
величини), де $\tau_E$~--- час утримання енергії. Жодна матеріальна стінка такої
температури не витримає, а~контакт зі стінкою охолодив би плазму. Тому плазму треба
утримувати без стінок. Є два підходи: інерційний (надвисока густина й~надкороткий час,
$\tau\sim10^{-10}$~с) і~магнітний ($n\sim10^{14}~\mathrm{см}^{-3}$, $\tau\sim1$~с), який
розглядається нижче.

Магнітне утримання ґрунтується на двох фактах, з~якими ми вже зустрічалися: заряджені
частинки прив'язані до силових ліній (розд.~\ref{sec:particle-motion}), а~плазма має
діамагнітні властивості, тому магнітне поле для неї~--- \enquote{подушка}. Ідею
запропонували А.~Д.~Сахаров та І.~Є.~Тамм у~1950~р.

%% --------------------------------------------------------

\section{Відкриті пастки: магнітне дзеркало}\index{Пастка!відкрита|textbf}\index{Магнітне дзеркало}
%% --------------------------------------------------------

Принцип дзеркала (умова відбиття~\eqref{eq:mirror-condition}, конус втрат, рис.~\ref{pic:magnetic-mirror})
розглянуто в~розд.~\ref{sec:magnetic-mirror}. Тут нагадаємо лише головний висновок для утримання:
однорідне поле не тримає частинку вздовж лінії, а~неоднорідність разом зі збереженням $\mu$
утримує лише ту частину розподілу, що лежить поза конусом втрат. Нагріті частинки
розсіюються кулонівськими зіткненнями в~конус втрат за час
$\tau_\mathrm{loss}\sim\nu_{ii}^{-1}\ln(B_{\max}/B_{\min})$ (розд.~\ref{sec:plasma-collisions}; темп втрат
задають іони, бо швидші електрони утримує амбіполярний потенціал), тому відкриті пастки \enquote{протікають} значно швидше, ніж закриті; для довгоживучої
плазми їх потрібно доповнювати іншими засобами (амбіполярні потенціали, тандемні пробки).
Крім того, дзеркальні конфігурації з~\enquote{поганою} кривиною (центр кривини силових ліній
лежить з~боку плазми) схильні до жолобкових нестійкостей (розд.~\ref{sec:rayleigh-taylor}); конфігурації min-$B$, у~яких кривина
\enquote{добра} (центр кривини з~боку вакууму), від них вільні.

%% --------------------------------------------------------

\section{Пінч-ефект}
%% --------------------------------------------------------

Плазму можна створити електричним розрядом. Далі виникає задача утримати плазмовий шнур
від розпливання, спричиненого газокінетичним тиском. Одним із~засобів є \emphx{пінч-ефект}\index{Пінч-ефект|textbf}
(від англ.~\emph{pinch}~--- стискання)~--- властивість струмового каналу зменшувати свій
перетин під дією магнітного поля, створеного власним струмом. Фізичну причину зрозуміти
легко: паралельні струми притягуються (сила Ампера на одиницю довжини
$dF/dl=2I_1I_2/c^2d$), і~нитки струму в~шнурі стискаються до осі.

Пінч-ефект описав Беннет 1934~р. стосовно потоків швидких заряджених частинок
у~газорозрядній плазмі, а~назву \enquote{пінч-ефект} 1937~р. запропонував Тонкс, який
вивчав дуговий розряд. Розрізняють два типи (рис.~\ref{pic:pinch-types}):
\begin{enumerate}
	\item $z$-пінч: струм тече вздовж шнура й~взаємодіє з~власним азимутальним полем;
	\item $\theta$-пінч: швидко зростаюче поздовжнє поле індукує в~плазмі азимутальні
	  (кільцеві) струми, що взаємодіють з~цим полем і~стискають шнур.
\end{enumerate}
Розглянемо рівновагу обох конфігурацій, а~потім їхню стійкість.

\begin{figure}[ht]\centering
	\localinput{pinch-types.tikz}
	\caption{Зліва~--- $z$-пінч, справа~--- $\theta$-пінч}
	\label{pic:pinch-types}
\end{figure}

%% --------------------------------------------------------
\section{Рівновага \texorpdfstring{$z$}{z}-пінча. Співвідношення Беннета}
%% --------------------------------------------------------

Розглянемо нескінченний циліндричний шнур радіуса $R$ із~поздовжнім струмом $I$
(рис.~\ref{pic:z-pinch-equilibrium}). Теорема про циркуляцію на колі радіуса $r$ дає
азимутальне поле
% ~~~~~~~~~~~~~~~~~~~~~~~~~~~~~~~~~~~~~~~~~~~~~~~~~~~~~~~~
\begin{equation*}
	2\pi rB_\theta=\frac{4\pi}cI(r)\quad\Rightarrow\quad B_\theta(r)=\frac{2I(r)}{cr},\qquad I(r)=\int_0^r2\pi r'j\,dr',
\end{equation*}
% ~~~~~~~~~~~~~~~~~~~~~~~~~~~~~~~~~~~~~~~~~~~~~~~~~~~~~~~~
а~умова рівноваги~\eqref{eq:equilibrium} у~радіальній проекції $dp/dr=-jB_\theta/c$,
оскільки $\vect j\times\vect B$ напрямлено до осі. Виразивши $j=\dfrac c{4\pi
r}\dfrac{d(rB_\theta)}{dr}$, знаходимо
% ~~~~~~~~~~~~~~~~~~~~~~~~~~~~~~~~~~~~~~~~~~~~~~~~~~~~~~~~
\begin{equation}\label{eq:z-pinch-general}
	\frac{d}{dr}\left(p+\frac{B_\theta^2}{8\pi}\right)=-\frac{B_\theta^2}{4\pi r}.
\end{equation}
% ~~~~~~~~~~~~~~~~~~~~~~~~~~~~~~~~~~~~~~~~~~~~~~~~~~~~~~~~~~~~~~~~
Права частина~--- це натяг азимутальних силових ліній (\enquote{обручевий натяг}), який
додатково стискає шнур. Тому сума газового та магнітного тисків $p+B_\theta^2/8\pi$
\emph{не} стала: вона спадає до периферії.

\begin{figure}[ht]\centering
	\localinput{z-pinch-equilibrium.tikz}
	\caption{До виведення умови рівноваги $z$-пінча. Контур $\Gamma$ показано штриховою
	лінією}
	\label{pic:z-pinch-equilibrium}
\end{figure}

\subsection*{Окремий випадок: однорідний струм}

Нехай $j=\mathrm{const}=I/\pi R^2$, тоді $B_\theta=2\pi jr/c$ і~\eqref{eq:z-pinch-general}
інтегрується: з~урахуванням $p(R)=0$
% ~~~~~~~~~~~~~~~~~~~~~~~~~~~~~~~~~~~~~~~~~~~~~~~~~~~~~~~~
\begin{equation*}
	p(r)=\frac{\pi j^2}{c^2}\left(R^2-r^2\right)=\frac{B_\theta^2(R)-B_\theta^2(r)}{4\pi}.
\end{equation*}
% ~~~~~~~~~~~~~~~~~~~~~~~~~~~~~~~~~~~~~~~~~~~~~~~~~~~~~~~~
Бачимо, що саме для $j=\mathrm{const}$ виконується $p+B_\theta^2/4\pi=\mathrm{const}$
(а~не $B_\theta^2/8\pi$: різницю дає обручевий натяг). Якщо концентрація всіх частинок $n$
стала, $p=nT$, то профіль температури
% ~~~~~~~~~~~~~~~~~~~~~~~~~~~~~~~~~~~~~~~~~~~~~~~~~~~~~~~~
\begin{equation*}
	T(r)=\frac{\pi}{nc^2}\left(\frac I{\pi R^2}\right)^2(R^2-r^2).
\end{equation*}
% ~~~~~~~~~~~~~~~~~~~~~~~~~~~~~~~~~~~~~~~~~~~~~~~~~~~~~~~~
Зокрема на осі $N_1T(0)=I^2/c^2$, де $N_1=\pi R^2n$~--- погонна густина
частинок.

\subsection*{Загальне співвідношення Беннета}\index{Беннета співвідношення|textbf}

Результат можна одержати без припущень про профіль. Помножимо $dp/dr=-jB_\theta/c$ на
$r^2$ й~проінтегруємо від $0$ до $R$. Ліва частина дає (інтегрування частинами, $p(R)=0$)
$-2\int_0^Rpr\,dr$, а~у~правій, використавши $j=\frac c{4\pi r}\frac{d(rB_\theta)}{dr}$,
маємо
% ~~~~~~~~~~~~~~~~~~~~~~~~~~~~~~~~~~~~~~~~~~~~~~~~~~~~~~~~
\begin{equation*}
	-\frac1c\int_0^Rr^2jB_\theta\,dr=-\frac1{8\pi}\int_0^R\frac{d\,(rB_\theta)^2}{dr}\,dr=-\frac{[RB_\theta(R)]^2}{8\pi}=-\frac{I^2}{2\pi c^2}.
\end{equation*}
% ~~~~~~~~~~~~~~~~~~~~~~~~~~~~~~~~~~~~~~~~~~~~~~~~~~~~~~~~
Звідси $\int_0^R2\pi p\,r\,dr=I^2/2c^2$, або
% ~~~~~~~~~~~~~~~~~~~~~~~~~~~~~~~~~~~~~~~~~~~~~~~~~~~~~~~~
\begin{equation}\label{eq:Bennett}
	\tcbhighmath{N_1\overline T=\frac{I^2}{2c^2}}.
\end{equation}
% ~~~~~~~~~~~~~~~~~~~~~~~~~~~~~~~~~~~~~~~~~~~~~~~~~~~~~~~~~~~~~~~
Тут $N_1$~--- повне число частинок (електронів і~іонів) на одиницю довжини шнура,
$\overline T$~--- усереднена за кількістю частинок температура. Це \emphx{співвідношення
Беннета}. Воно показує, що для утримання шнура із~заданими $N_1$ і~$\overline T$
\emph{потрібен цілком певний струм}, незалежно від радіуса й~розподілу струму.
У~випадку~$n=\mathrm{const}$, $j=\mathrm{const}$ маємо $\overline T=T(0)/2$,
і~\eqref{eq:Bennett} збігається з~отриманим вище. Наприклад, шнур з~концентрацією
електронів й~іонів по $10^{16}~\mathrm{см}^{-3}$ при радіусі 1~см, нагрітий в~середньому
до 1~кеВ, потребує струму $I\approx1{,}4\times10^5$~А (аналогічна задача~\ref{prb:PlasmaBennett}).

%% --------------------------------------------------------
\section{Рівновага \texorpdfstring{$\theta$}{θ}-пінча}\label{sec:theta-pinch}\index{Тета-пінч|textbf}
%% --------------------------------------------------------

У~$\theta$-пінчі (рис.~\ref{pic:pinch-types}, справа) плазму оточує одновиткова котушка,
через яку розряджають конденсаторну батарею. Поздовжнє поле $B_z$ за мікросекунди зростає
до десятків кілогаусів. Вихрове електричне поле $E_\theta$ наводить у~поверхневому шарі
плазми кільцевий струм $j_\theta$. Його напрям такий, що власне поле струму зменшує
зовнішнє поле всередині шнура: плазма поводиться як діамагнетик
(розд.~\ref{sec:diamagnetism}). Сила $\vect j\times\vect B/c$ напрямлена до осі, і~поле
працює як \emphx{магнітний поршень}, що стискає шнур.

У~рівновазі радіальна проекція~\eqref{eq:equilibrium} має вигляд $dp/dr=j_\theta B_z/c$.
Підставивши $j_\theta=-\dfrac c{4\pi}\dfrac{dB_z}{dr}$, одержуємо
% ~~~~~~~~~~~~~~~~~~~~~~~~~~~~~~~~~~~~~~~~~~~~~~~~~~~~~~~~
\begin{equation}\label{eq:theta-pinch}
	\frac{d}{dr}\left(p+\frac{B_z^2}{8\pi}\right)=0\quad\Rightarrow\quad p(r)+\frac{B_z^2(r)}{8\pi}=\frac{B_0^2}{8\pi},
\end{equation}
% ~~~~~~~~~~~~~~~~~~~~~~~~~~~~~~~~~~~~~~~~~~~~~~~~~~~~~~~~
де $B_0$~--- поле між плазмою та котушкою, де $p=0$. Це просто баланс
тисків~\eqref{eq:pressure-balance}. Порівняно з~$z$-пінчем~\eqref{eq:z-pinch-general}
у~правій частині немає обручевого натягу: силові лінії поля $B_z$ прямі. Тому сума газового
та магнітного тисків стала по перерізу.

З~\eqref{eq:theta-pinch} випливає, що $\beta=8\pi p/B_0^2=1-B_z^2/B_0^2\leqslant1$.
Межа $\beta=1$ досягається, коли поле повністю витіснене з~осі. На відміну від токамака,
де $\beta$ становить кілька відсотків (розд.~\ref{sec:beta-limits}), у~$\theta$-пінчах
одержують $\beta\approx0{,}7$--$0{,}9$. Отже, $\theta$-пінч використовує магнітне поле
найефективніше. Максимальний тиск задає лише поле котушки. Наприклад, при
$B_0=50$~кГс маємо $p_{\max}=B_0^2/8\pi\approx10^8~\mathrm{дин/см^2}\approx100$~атм.
При $T_e=T_i=1$~кеВ це відповідає $n=p/2T\approx3\times10^{16}~\mathrm{см}^{-3}$.

\emph{Нагрівання.} Плазма нагрівається у~два етапи. Спочатку швидкий фронт поля (час
наростання $\lesssim1$~мкс) збуджує стискну ударну хвилю. Вона за один прохід нагріває
іони до сотень електронвольт. Далі поле зростає повільніше, і~відбувається адіабатичне
стискання вмороженої плазми. Саме на такій установці, \emph{Scylla} (Лос-Аламос, 1958~р.),
уперше надійно зареєстрували термоядерні нейтрони з~лабораторної плазми.

\emph{Стійкість і~втрати.} Силові лінії прямі, а~поле на межі плазми монотонно зростає
назовні. Тому лінійний $\theta$-пінч стійкий щодо перетяжок і~викривлень, які руйнують
$z$-пінч (розд.~\ref{sec:pinch-instabilities}). Проте вздовж поля ніщо не утримує частинки,
і~плазма вільно витікає через торці. Час утримання визначається часом прольоту іона на
половину довжини, $\tau\sim L/2v_{Ti}$. Для дейтерієвої плазми з~$T=1$~кеВ
($v_{Ti}\approx2\times10^7$~см/с) при $L=1$~м маємо $\tau\approx2$~мкс і~$n\tau\sim10^{11}~\mathrm{см}^{-3}\,\mathrm{с}$.
Це на три порядки менше за критерій Лоусона~\eqref{eq:lawson}. Закрити торці, згорнувши
шнур у~тор, теж не вдалося. Як і~для будь-якого чисто тороїдального поля
(розд.~\ref{sec:toroidal}), тороїдальний $\theta$-пінч (\emph{Scyllac}, 1970-ті) не має
рівноваги без додаткових гвинтових полів. Його нащадок~--- \emphx{конфігурація з~оберненим
полем} (FRC). У~ній на початку розряду всередині шнура вморожують зворотне поле, і~силові
лінії замикаються самі на себе в~плазмі.

%% --------------------------------------------------------
\section{Нестійкості пінча}\label{sec:pinch-instabilities}\index{Пінч!нестійкості}
%% --------------------------------------------------------

Рівновага ще не гарантує існування. Плазма~--- об'єкт, у~якому розвиваються різні
нестійкості, що порушують її однорідність і~форму та обмежують час життя. Для
термоядерного синтезу, де плазму потрібно утримувати довго, це критично.

\subsection*{Перетяжки (\texorpdfstring{$m=0$}{m = 0})}

Нехай випадково на шнурі виникло звуження (рис.~\ref{pic:pinch-instability}). Струм $I$
сталий, тому густина струму $j=I/\pi R^2$ в~перетяжці зростає, а~разом із~нею зростає поле
на поверхні $B_\theta(R)=2I/cR$. Тиск цього поля $B_\theta^2/8\pi=I^2/2\pi c^2R^2\propto
R^{-2}$ збільшується, а~сила Ампера на одиницю об'єму $F_A=jB_\theta/c=2I^2/(\pi c^2R^3)$
ще швидше. Тиск плазми в~перетяжці майже не підвищується (плазма витискається вздовж осі).
Тому нескомпенсована сила стискання зростає, звуження поглиблюється~--- розвивається
нестійкість. Для лабораторного шнура ($n\sim10^{16}~\mathrm{см}^{-3}$, $R\sim1$~см)
характерний час росту $\sim R/v_{A\theta}\sim10^{-8}$~с, де
$v_{A\theta}=B_\theta/\sqrt{4\pi\rho_m}$. Для прикладу Беннета ($B_\theta\approx2{,}8\times10^{4}$~Гс, $n_i=10^{16}~\mathrm{см}^{-3}$, дейтерій)
$v_{A\theta}\approx4\times10^{7}$~см/с, $R/v_{A\theta}\approx2{,}5\times10^{-8}$~с.

\subsection*{Викривлення (\texorpdfstring{$m=1$}{m = 1})}

Якщо шнур випадково прогнувся, то з~увігнутого (внутрішнього) боку вигину кільцеві силові лінії $B_\theta$ згущуються,
а~з~опуклого розріджуються. Магнітний тиск $B^2/8\pi$ з~увігнутого боку більший, він штовхає шнур ще далі в~бік опуклості,
і~вигин наростає: що сильніше прогин, то більша різниця тисків, тобто нестійкість.

\begin{figure}[ht]\centering
	\localinput{pinch-instability.tikz}
	\caption{Нестійкості $z$-пінча: перетяжки та викривлення шнура}
	\label{pic:pinch-instability}
\end{figure}

\subsection*{Стабілізація поздовжнім полем}\index{Стабілізація пінча поздовжнім полем}

Щоб боротися з~перетяжками, в~плазму \enquote{вморожують} поздовжнє магнітне поле $B_z$.
У~ідеально провідній плазмі магнітний потік через поперечний переріз зберігається: $\pi
R^2B_z=\mathrm{const}$, тому при звуженні шнура $B_z\propto R^{-2}$, а~відповідний тиск
$B_z^2/8\pi\propto R^{-4}$ зростає \emph{швидше} за тиск азимутального поля ($\propto
R^{-2}$). Вимагатимемо рівноваги на поверхні $p+B_z^2/8\pi=B_\theta^2(R)/8\pi$
й~порівняємо зміни при малому стисканні $\delta R$ (тиск $p$ не змінюється):
% ~~~~~~~~~~~~~~~~~~~~~~~~~~~~~~~~~~~~~~~~~~~~~~~~~~~~~~~~
\begin{equation*}
	\delta\frac{B_z^2}{8\pi}=-4\,\frac{B_z^2}{8\pi}\frac{\delta R}R,\qquad
	\delta\frac{B_\theta^2}{8\pi}=-2\,\frac{B_\theta^2}{8\pi}\frac{\delta R}R.
\end{equation*}
% ~~~~~~~~~~~~~~~~~~~~~~~~~~~~~~~~~~~~~~~~~~~~~~~~~~~~~~~~
Внутрішній тиск при стисканні зростає більше за зовнішній, якщо $4B_z^2>2B_\theta^2$,
тобто $B_z^2>B_\theta^2/2$. Тоді виникає повертаюча сила~--- перетяжки стійкі. (Це оцінка
в~спрощеній моделі, але вона правильно передає суть: чим більше поздовжнє поле, тим
стійкіший шнур; у~цій моделі тиск плазми вважається незмінним, а~$B_z$~--- \enquote{вмороженим} у~шнур і~стисненим разом із~ним.)

Проти викривлення потрібна ще сильніша стабілізація. Умова Крускала--Шафранова для
шнура довжиною $L$ й~радіусом $a$ вимагає
% ~~~~~~~~~~~~~~~~~~~~~~~~~~~~~~~~~~~~~~~~~~~~~~~~~~~~~~~~
\begin{equation*}
	q=\frac{2\pi a}L\,\frac{B_z}{B_\theta}>1.
\end{equation*}
% ~~~~~~~~~~~~~~~~~~~~~~~~~~~~~~~~~~~~~~~~~~~~~~~~~~~~~~~~~
Тобто силові лінії мають зробити більше одного оберту навколо осі на довжині шнура.
У~стабілізованому пінчі $B_z$ значно перевищує $B_\theta$, а~рівновага визначається не
лише струмом, і~ми поступово переходимо до сучасних систем. Крім перетяжок і~викривлень
існують інші нестійкості (перестановочні, балонні, розривні моди~--- останні пов'язані з~перез'єднанням, розд.~\ref{sec:reconnection}), які визначають межі
параметрів реальних установок.

%% --------------------------------------------------------
\section{Тороїдальні системи: токамак і~стеларатор}\label{sec:toroidal}
%% --------------------------------------------------------

Відкриті пастки (дзеркала, лінійні пінчі) втрачають частинки через торці. Інший спосіб~---
замкнути силові лінії в~тор. Однак чисто тороїдальне поле $B_\varphi\propto1/R$
неоднорідне, тому виникають градієнтний і~відцентровий дрейфи~\eqref{eq:gradB-drift}
з~вертикальною швидкістю
% ~~~~~~~~~~~~~~~~~~~~~~~~~~~~~~~~~~~~~~~~~~~~~~~~~~~~~~~~
\begin{equation*}
	v_d\approx\frac{2c\,T}{|q|BR},
\end{equation*}
% ~~~~~~~~~~~~~~~~~~~~~~~~~~~~~~~~~~~~~~~~~~~~~~~~~~~~~~~~
протилежною для іонів і~електронів. Для $T=10$~кеВ, $B=5\times10^4$~Гс,
$R=6\times10^2$~см маємо $v_d\approx7\times10^4$~см/с, тобто за секунду частинка дрейфувала б
на сотні метрів. Дрейф розділяє заряди, виникає вертикальне електричне поле, і~дрейф
у~схрещених полях~\eqref{eq:ExB} виносить уже \emph{обидва} сорти частинок назовні на
стінку. Отже, чисто тороїдальне поле не утримує.

Вихід~--- закрутити силові лінії, додавши до тороїдального поля полоїдальну складову
$B_\theta$. Тоді лінії гвинтоподібно намотуються на вкладені магнітні поверхні,
і~частинка, рухаючись уздовж лінії, по черзі опиняється вгорі й~внизу тора: вертикальні
дрейфи усереднюються, а~заряди закорочуються вздовж лінії. Закрученість характеризують
\emphx{запасом стійкості}\index{Запас стійкості|textbf} $q=rB_\varphi/(RB_\theta)$~--- локальним співвідношенням
тороїдального та полоїдального полів. Фізично $q$ показує, скільки разів силова лінія обходить тор вздовж (тороїдально) за один оберт навколо осі шнура (полоїдально);
назва \enquote{запас стійкості}~--- це запас до порогу $q=1$. Для шнура довжиною $L=2\pi R$ маємо $q=2\pi aB_z/(LB_\theta)=aB_\varphi/(RB_\theta)$~---
той самий $q$ в~обох означеннях. Важливий також \emph{шир} $s=(r/q)\,dq/dr$: на сусідніх магнітних поверхнях лінії мають різний крок гвинта,
тому збурення, витягнуте вздовж лінії на одній поверхні, на сусідній \enquote{розрізається}; шир стабілізує. У~простій циліндричній моделі критерій Крускала~---
Шафранова для основної $m=1$-моди приводить до умови порядку $q(a)>1$. Для реального
токамака умови МГД-стійкості складніші й~визначаються профілями струму, тиску та
геометрією магнітних поверхонь.

Закрутку створюють двома способами (рис.~\ref{pic:toroidal-trap}). У~\emphx{токамаку}\index{Токамак|textbf}
(скорочення від рос.~\emph{тороидальная камера с магнитными катушками}) тороїдальне поле
створюють зовнішні котушки, а~полоїдальне~--- струм, що тече самою плазмою (індукований
трансформатором). Токамак розвинули в~СРСР з~ідей Тамма й~Сахарова; результати установки
Т-3 (1968~р., під керівництвом Л.~А.~Арцимовича) стали поворотним пунктом, після якого
токамаки почали будувати в~усьому світі. У~\emphx{стеларатор}\index{Стелларатор|textbf} (Спітцер, 1951~р.) закрутку
створюють зовнішні гвинтові котушки, а~струму в~плазмі немає, що дозволяє працювати
в~стаціонарному режимі, але вимагає складної тривимірної геометрії котушок. Найбільший
токамак, що будується,~--- ITER ($R\approx6{,}2$~м, $a\approx2$~м, $B\approx5{,}3$~Тл,
$I=15$~МА); саме такі параметри ми використовуватимемо в~оцінках нижче.

\begin{figure}[ht]\centering
	\localinput{toroidal.tikz}
	\caption{Принцип тороїдального магнітного утримання: силові лінії замкнені
	й~закручені, лежать на вкладених магнітних поверхнях}
	\label{pic:toroidal-trap}
\end{figure}

%% --------------------------------------------------------
\section{Захоплені частинки й~неокласичний перенос}\label{sec:neoclassical}\index{Неокласичний перенос|textbf}\index{Захоплені частинки|textbf}
%% --------------------------------------------------------

У~токамаку $B_\varphi\propto1/R$: поле сильніше на внутрішній стороні тора ($B_{\max}$) і~слабше
на зовнішній ($B_{\min}$). Отже, тор містить магнітне дзеркало (розд.~\ref{sec:magnetic-mirror}), розтягнуте по
периметру. Умова відбиття $\sin^2\theta\ge B_{\min}/B_{\max}=(1-\varepsilon)/(1+\varepsilon)$, де
$\varepsilon=r/R$~--- обернене аспектне відношення, означає (кути беруться на зовнішній стороні, де $B=B_{\min}$)
$\sin^2\theta\ge\dfrac{1-\varepsilon}{1+\varepsilon}\Rightarrow\cos^2\theta\le\dfrac{2\varepsilon}{1+\varepsilon}\Rightarrow
v_\parallel/v\lesssim\sqrt{2\varepsilon}$: такі частинки \emph{захоплені} на зовнішній стороні й~не оббігають
поперечний переріз повністю. Їхня частка $f_t\approx\sqrt{2\varepsilon}$: для ізотропного розподілу частка
частинок у~конусі $|\cos\theta|<x$ дорівнює $x$. Для $\varepsilon=0{,}3$ маємо $\sqrt{0{,}6}\approx0{,}77\sim0{,}8$~---
це велика частка, а~не мала поправка (точніше, $f_t\approx\sqrt{2\varepsilon/(1+\varepsilon)}\approx0{,}68$). Інші частинки \emph{пролітні}.

\begin{figure}[ht]\centering
	\localinput{banana.tikz}
	\caption{Полоїдальний переріз токамака: пролітна частинка оббігає магнітну поверхню, а~захоплена
	описує \enquote{бананову} орбіту навколо екваторіальної площини на зовнішній стороні}
	\label{pic:banana}
\end{figure}

У~проєкції на полоїдальну площину траєкторія захопленої частинки має форму банана
(рис.~\ref{pic:banana}). Оцінимо її ширину: поздовжня швидкість захопленої частинки $v_\parallel\sim\sqrt\varepsilon\,v$;
час півоберту по банану $\tau_b\sim qR/v_\parallel\sim qR/(\sqrt\varepsilon\,v)$ (множник $q$ тому, що довжина силової лінії за один
полоїдальний оберт дорівнює $2\pi qR$); тороїдальний дрейф $v_d\sim v\,r_L/R$. Отже, $\Delta_b\sim v_d\tau_b\sim q\,r_L/\sqrt\varepsilon$:
радіальна ширина за порядком величини
% ~~~~~~~~~~~~~~~~~~~~~~~~~~~~~~~~~~~~~~~~~~~~~~~~~~~~~~~~
\begin{equation}\label{eq:banana-width}
	\Delta_b\sim\frac{q\,r_L}{\sqrt\varepsilon},
\end{equation}
% ~~~~~~~~~~~~~~~~~~~~~~~~~~~~~~~~~~~~~~~~~~~~~~~~~~~~~~~~
де $q$~--- запас стійкості, $r_L$~--- ларморівський радіус у~повному полі. Це кратно
більше за $r_L$ (для $q=2$, $\varepsilon=0{,}3$ у~$\sim4$ рази). Кулонівські
зіткнення зсувають частинку на $\Delta_b$ вбік; щоб вона покинула банан, достатньо розсіяти її по
куту на $\Delta\theta\sim\sqrt\varepsilon$. Кут розсіюється \emph{дифузійно}, $(\Delta\theta)^2\sim\nu t$, тож щоб змінити його
на $\sqrt\varepsilon$, потрібен час $t\sim\varepsilon/\nu$, тобто ефективна частота розсіяння $\nu_\mathrm{eff}\sim\nu/\varepsilon$
(вона більша за $\nu$, бо для виходу з~банана досить малого кутового зсуву).
Оцінимо дифузію (крок $\Delta_b$, частота кроків $\nu_\mathrm{eff}$, частка частинок $f_t$):
% ~~~~~~~~~~~~~~~~~~~~~~~~~~~~~~~~~~~~~~~~~~~~~~~~~~~~~~~~
\begin{equation}\label{eq:neoclassical}
	D_\mathrm{neo}\sim f_t\,\Delta_b^2\,\nu_\mathrm{eff}\sim q^2\,\varepsilon^{-3/2}\,r_L^2\nu.
\end{equation}
% ~~~~~~~~~~~~~~~~~~~~~~~~~~~~~~~~~~~~~~~~~~~~~~~~~~~~~~~~
Порівняно з~класичною дифузією $D_\perp\sim r_L^2\nu$ перенос збільшується в~$q^2\varepsilon^{-3/2}$
разів, тобто в~$\sim20$--$50$ разів у~типових токамаках. Це \emph{неокласичний} перенос
(режим бананів, поки $\nu\,qR/(\varepsilon^{3/2}v_T)<1$). Він є результатом тороїдальної геометрії,
не залежить від турбулентності, але все одно нижчий за спостережуваний у~електронному каналі~---
різницю (див.~розд.~\ref{sec:drift-waves}) дає аномальний перенос.

Захоплені частинки мають ще один наслідок: у~плазмі з~градієнтом тиску виникає
\emph{бутстреп-струм} $j_\mathrm{bs}\sim-\sqrt\varepsilon\,(c/B_\theta)\,dp/dr$, який створюється самою плазмою
без трансформатора й~дозволяє токамаку працювати в~стаціонарному режимі.

\emph{Звідки він береться.} Захоплені частинки на бананах рухаються в~протилежних напрямках уздовж внутрішньої та зовнішньої \enquote{ніжок}
банана. Якщо є градієнт густини, на даному радіусі банани, що приходять із~внутрішнього боку, густіші за ті, що з~зовнішнього,
і~виникає нескомпенсований струм (аналог діамагнітного струму намагнічування, але від бананів). Через тертя він передається пролітним частинкам,
звідки й~оцінка вище. Практично: у~стаціонарних сценаріях ITER і~в~проєктах DEMO розраховують на частку
бутстреп-струму $\sim50\%$ і~більше (в~індуктивному базовому сценарії ITER~--- лише $\sim15$--$20\%$),
що наближає токамак до стаціонарної роботи без індукції.

%% --------------------------------------------------------
\section{Рівновага тора: зсув Шафранова та межа за тиском}\label{sec:beta-limits}
%% --------------------------------------------------------

Плазмовий тор, як і~надувний обруч, прагне розширитися: газовий тиск і~\enquote{обручева} сила струму
витискають кільце назовні. Для рівноваги потрібне вертикальне поле $B_v$, а~самі магнітні поверхні
зсуваються вбік зовнішньої стінки відносно центра камери~--- \emphx{зсув Шафранова}\index{Шафранова зсув|textbf}
$\Delta\sim(a^2/R)(\beta_p+l_i/2)$, де $\beta_p=8\pi\langle p\rangle/B_\theta^2$~--- полоїдальна бета,
$l_i$~--- внутрішня індуктивність. Чим більший тиск плазми, тим більший зсув; при $\beta_p\sim R/a$ рівновага
взагалі зникає.

Є й~жорсткіші межі, що пов'язані зі~стійкістю (розд.~\ref{sec:instability}). Поздовжній струм обмежений
кінк-нестійкістю: $q(a)>1$, а~практично $q(a)\gtrsim2$--$3$. Тиск обмежений балонними й~кінк-модами; для них
емпірично встановлено \emphx{межу Тройона}\index{Тройона межа|textbf}
% ~~~~~~~~~~~~~~~~~~~~~~~~~~~~~~~~~~~~~~~~~~~~~~~~~~~~~~~~
\begin{equation}\label{eq:troyon}
	\tcbhighmath{\beta_{\max}[\%]=\beta_N\,\frac{I[\text{МА}]}{a[\text{м}]\,B[\text{Тл}]},\qquad\beta_N\approx2{,}8,}
\end{equation}
% ~~~~~~~~~~~~~~~~~~~~~~~~~~~~~~~~~~~~~~~~~~~~~~~~~~~~~~~~
тобто $\beta\sim$ кілька відсотків. Для $I=15$~МА, $a=2$~м, $B=5$~Тл маємо
$\beta_{\max}\approx4\%$. Оскільки $P_f\propto\beta^2B^4$ (розд.~\ref{sec:mhd}), саме межа Тройона й~максимальне
поле надпровідних магнітів визначають розмір електростанції. Окрім того, густина обмежена \emph{межею Грінвальда}
$n_G[10^{20}\,\text{м}^{-3}]=I[\text{МА}]/(\pi a^2[\text{м}^2])\approx1{,}2$ для цих параметрів
($\approx1{,}2\times10^{14}~\mathrm{см}^{-3}$)~--- не випадково саме такі густини ми розглядали у~попередніх розділах.

\begin{Attention}
	Бета порівнює тиски, але сама по~собі не~характеризує стійкість: утримання з~малим $\beta$ може бути нестійким (наприклад, при $q(a)<1$), а~межа Тройона~--- емпіричне,
	а~не~точне співвідношення.
\end{Attention}

%% --------------------------------------------------------
\section{Як це вимірюють: магнітна реконструкція рівноваги}\label{sec:equilibrium-diagnostics}\index{Діагностика плазми!магнітна реконструкція рівноваги}
%% --------------------------------------------------------

Форму й~положення плазмового шнура не можна побачити всередині, але їх можна відновити із~зовнішніх магнітних вимірів. Пояс Роговського дає повний струм~$I$,
петлі потоку й~зонди~$B_\theta$ вздовж стінки (розд.~\ref{sec:magnetic-diagnostics}) дають полоїдальне поле на~контурі. Поза плазмою поле безвихрове, тож за~значеннями на~контурі
розв'язують рівняння Града--Шафранова й~знаходять магнітні поверхні, положення осі й~зсув Шафранова~$\Delta$. Зворотно, виміряний~$\Delta$ задає комбінацію
% ~~~~~~~~~~~~~~~~~~~~~~~~~~~~~~~~~~~~~~~~~~~~~~~~~~~~~~~~
\begin{equation}\label{eq:shafranov-diagnostic}
	\tcbhighmath{\beta_p+\frac{l_i}{2}\sim\frac{R}{a^2}\,\Delta,}
\end{equation}
% ~~~~~~~~~~~~~~~~~~~~~~~~~~~~~~~~~~~~~~~~~~~~~~~~~~~~~~~~
а~разом з~діамагнітним вимірюванням енергії (розд.~\ref{sec:magnetic-diagnostics}) розділяє внесок тиску~$\beta_p$ і~внутрішньої індуктивності~$l_i$ (профілю струму).
\begin{Estimate}
Для $a=0{,}5$~м, $R=1{,}5$~м, $\Delta=10$~см: $\beta_p+l_i/2\approx0{,}6$.
\end{Estimate}

\emph{Межі застосовності.} Зовнішні виміри однозначно задають лише зовнішні поверхні; внутрішній профіль струму $q(r)$ потрібно знаходити з~додаткових даних
(поляриметрія, розд.~\ref{sec:faraday-diagnostics}, або штарківське розщеплення ліній інжектованого пучка), а~у~реальному часі реконструкцію виконують спрощеними моделями.

%% --------------------------------------------------------
\begin{Summary}
\begin{itemize}
\item Магнітне поле обмежує рух поперек силових ліній, а~вздовж них~--- ні, тому геометрія поля принципово важлива.
\item Відкриті пастки (магнітне дзеркало) мають конус втрат; замкнені конфігурації вимагають закрутки силових ліній.
\item Рівновага $z$-пінча встановлюється балансом газокінетичного тиску, магнітного тиску й~натягу азимутальних силових ліній (обручевий натяг); такий пінч нестійкий до перетяжок і~викривлень.
\item У~$\theta$-пінчі силові лінії прямі, тож $p+B_z^2/8\pi=\mathrm{const}$, $\beta$ близьке до одиниці, а~конфігурація МГД-стійка; головна вада~--- швидке витікання плазми через торці.
\item Токамак і~стеларатор утримують плазму завдяки гвинтовому полю; захоплені частинки призводять до неокласичного переносу.
\item Тиск плазми обмежений рівновагою та стійкістю (зсув Шафранова, межа за $\beta$).
\item Положення й~форму шнура, зсув Шафранова та комбінацію $\beta_p+l_i/2$ відновлюють із~зовнішніх магнітних вимірів (розд.~\ref{sec:equilibrium-diagnostics}).
\end{itemize}
\end{Summary}

\Needspace{9\baselineskip}
\section*{Задачі}
%% --------------------------------------------------------

\begin{problem}\label{prb:BananaWidth}\lvl{2}
	Для $\varepsilon=a/R=0{,}3$ ($a=2$~м), $q=2$ знайдіть частку захоплених частинок $f_t$ і~ширину банана
	для іона дейтерію ($T=10$~кеВ, $B=5\times10^4$~Гс). Порівняйте з~$a$ і~оцініть
	$D_\mathrm{neo}/D_\mathrm{cl}$.
	\begin{solution}
		$f_t\approx\sqrt{2\varepsilon}\approx0{,}77$; $r_L\approx0{,}3$~см, $\Delta_b\sim qr_L/\sqrt\varepsilon\approx1$~см $\ll a=200$~см;
	$D_\mathrm{neo}/D_\mathrm{cl}\sim q^2\varepsilon^{-3/2}\approx24$.
	\end{solution}
\end{problem}

\begin{problem}\label{prb:Troyon}\lvl{2}
	Знайдіть межу Тройона й~межу Грінвальда для $I=2$~МА, $a=0{,}5$~м, $B=2$~Тл. Яким буде
	тиск плазми при досягненні цієї межі $\beta$?
	\begin{solution}
		$\beta_{\max}\approx2{,}8\cdot2/(0{,}5\cdot2)\approx5{,}6\%$; $n_G\approx2{,}5\times10^{20}~\mathrm{м}^{-3}=2{,}5\times10^{14}~\mathrm{см}^{-3}$;
	$p=\beta B^2/8\pi\approx9\times10^5~\mathrm{дин/см^2}\approx0{,}9$~атм.
	\end{solution}
\end{problem}

\begin{problem}\label{prb:PlasmaBennett}\lvl{2}
	За співвідношенням Беннета знайдіть струм, потрібний для утримання шнура радіуса
	$R=0{,}5$~см дейтерієвої плазми з~концентрацією електронів та іонів по $5\times10^{17}~\mathrm{см}^{-3}$ і~середньою
	температурою $200$~еВ. Яке при цьому поле $B_\theta$ на поверхні шнура? Оцініть час
	наростання перетяжки ($m=0$).
	\begin{solution}
		$N_1=2n\pi R^2\approx7{,}9\times10^{17}~\mathrm{см}^{-1}$, $I\approx2{,}2\times10^{5}$~А,
	$B_\theta=2I/cR\approx9\times10^{4}$~Гс; $\tau\sim R/v_A\sim2\times10^{-8}$~с.
	\end{solution}
\end{problem}

\begin{problem}\label{prb:ThetaPinch}\lvl{2}
	У~$\theta$-пінчі поле котушки $B_0=40$~кГс, а~поле на осі шнура $B_z(0)=20$~кГс.
	Знайдіть $\beta$ на осі, тиск плазми й~концентрацію дейтерієвої плазми при
	$T_e=T_i=1$~кеВ. Оцініть час витікання через торці для шнура довжиною $L=3$~м
	і~порівняйте $n\tau$ з~критерієм Лоусона~\eqref{eq:lawson}.
	\begin{solution}
		$\beta=1-B_z^2(0)/B_0^2=0{,}75$; $p=(B_0^2-B_z^2)/8\pi\approx4{,}8\times10^7~\mathrm{дин/см^2}\approx47$~атм;
	$n=p/2T\approx1{,}5\times10^{16}~\mathrm{см}^{-3}$; $v_{Ti}\approx2{,}2\times10^7$~см/с, $\tau\sim L/2v_{Ti}\approx7$~мкс;
	$n\tau\approx10^{11}~\mathrm{см}^{-3}\,\mathrm{с}$, на три порядки менше за потрібне.
	\end{solution}
\end{problem}

\begin{problem}\label{prb:SafetyFactor}\lvl{1}
	Оцініть запас стійкості $q(a)=aB_\varphi/(RB_\theta)$ для $a=0{,}5$~м, $R=1{,}5$~м, $B_\varphi=2$~Тл,
	$I=1$~МА, де $B_\theta=2I/(ca)$ у~СГС ($0{,}4$~Тл). Чи виконується умова Крускала--Шафранова $q(a)>1$?
	\begin{solution}
		$q\approx1{,}7$.
	\end{solution}
\end{problem}

\begin{problem}\label{prb:LossConeFraction}\lvl{2}
	Для пробкового відношення $R_m=B_{\max}/B_{\min}=5$ знайдіть частку частинок ізотропного розподілу,
	що потрапляють у~конус втрат. Як вона зміниться при $R_m=10$?
	\begin{solution}
		$1-\sqrt{1-1/R_m}\approx0{,}11$ і~$\approx0{,}05$.
	\end{solution}
\end{problem}

\begin{problem}\label{prb:ShafranovShift}\lvl{2}
	У~токамаку з~$a=0{,}5$~м, $R=1{,}5$~м, $B_\theta=0{,}4$~Тл середня концентрація $n_e=n_i=5\times10^{13}~\mathrm{см}^{-3}$, $T_e=T_i=1$~кеВ.
	Знайдіть $\beta_p=8\pi\langle p\rangle/B_\theta^2$ і~зсув Шафранова $\Delta$ за~\eqref{eq:shafranov-diagnostic} при $l_i=1$.
	\begin{solution}
		$\langle p\rangle=2nT\approx1{,}6\times10^5$~дин/см$^2$, $\beta_p\approx0{,}25$; $\Delta\approx(a^2/R)(\beta_p+l_i/2)\approx12$~см.
	\end{solution}
\end{problem}

\Needspace{6\baselineskip}
\begin{Questions}
\begin{quizlist}
\item Як працює магнітне дзеркало? Що таке конус втрат?

\item Чому $z$-пінч утримується? Чому в~загальному випадку
		  $p+B^2/8\pi\neq\mathrm{const}$?

\item Що виражає співвідношення Беннета? Що з~нього випливає для струму?

\item Чому в~$\theta$-пінчі $p+B_z^2/8\pi=\mathrm{const}$, а~в~$z$-пінчі ні? Чому
		  лінійний $\theta$-пінч стійкий, але не придатний для реактора?

\item Чому чисто тороїдальне поле не утримує плазму? Як це виправляють у~токамаку
		  та стелараторі?

\item Які частинки в~токамаку називають захопленими? Яку форму має їхня орбіта і~чому
		  неокласичний перенос перевищує класичний?

\item Які обмеження за струмом, тиском і~густиною існують у~токамаку?

\item Що таке запас стійкості $q$ і~навіщо потрібна умова $q(a)>1$?

\item Чому виникає бутстреп-струм?

\item Чому тор прагне розширитися і~як його утримують?

\end{quizlist}
\end{Questions}

\Closesolutionfile{answer}
